\section{Формула полной вероятности и формула Байеса}

Во многих задачах случайный эксперимент протекает в два этапа. На первом этапе реализуется одно из нескольких возможных условий (гипотез), а на втором — происходит интересующее нас событие.

\subsection{Полная группа гипотез}

\begin{definition}
События $H_1, H_2, \dots, H_n$ образуют \textbf{полную группу гипотез}, если выполняются три условия:
\begin{enumerate}
\item Они попарно несовместны: $H_i \cap H_j = \varnothing$ для любых $i \neq j$.
\item В результате эксперимента обязательно произойдет хотя бы одно из них: $\bigcup_{i=1}^n H_i = \Omega$.
\item Сумма их вероятностей равна единице: $\sum_{i=1}^n P(H_i) = 1$.
\end{enumerate}
\end{definition}

\subsection{Формула полной вероятности}

\begin{theorem}[Формула полной вероятности]
Вероятность события $A$, которое может наступить лишь при условии появления одного из несовместных событий $H_1, \dots, H_n$, образующих полную группу, равна сумме произведений вероятностей каждой из гипотез на соответствующую условную вероятность события $A$:
$$P(A) = \sum_{i=1}^n P(H_i) \cdot P(A\vert{}H_i).$$
\end{theorem}

\begin{proof}
Так как гипотезы $H_1, \dots, H_n$ образуют полное пространство $\Omega$, событие $A$ можно представить как его пересечение со всем пространством:
$$A = A \cap \Omega = (A \cap H_1) \cup (A \cap H_2) \cup \dots \cup (A \cap H_n).$$
События $(A \cap H_i)$ попарно несовместны. Применим аксиому сложения вероятностей:
$$P(A) = \sum_{i=1}^n P(A \cap H_i).$$
Согласно теореме умножения вероятностей $P(A \cap H_i) = P(H_i) \cdot P(A\vert{}H_i)$. Подставляем в сумму и получаем искомую формулу.
\end{proof}

\subsection{Формула Байеса (Теорема гипотез)}

\begin{theorem}[Формула Байеса]
Если событие $A$ уже произошло ($P(A) > 0$), то вероятность того, что оно было вызвано именно гипотезой $H_k$, пересчитывается (уточняется) по формуле:
$$P(H_k\vert{}A) = \frac{P(H_k) \cdot P(A\vert{}H_k)}{P(A)} = \frac{P(H_k) \cdot P(A\vert{}H_k)}{\sum_{i=1}^n P(H_i) \cdot P(A\vert{}H_i)}.$$
\end{theorem}

\begin{proof}
Запишем вероятность совместного наступления гипотезы $H_k$ и события $A$ двумя способами:
$$P(H_k \cap A) = P(A) \cdot P(H_k\vert{}A) \quad \text{и} \quad P(H_k \cap A) = P(H_k) \cdot P(A\vert{}H_k).$$
Приравняем правые части и разделим на $P(A)$:
$$P(H_k\vert{}A) = \frac{P(H_k) \cdot P(A\vert{}H_k)}{P(A)}.$$
\end{proof}

\subsection{Визуализация: Деревья вероятностей}

Для наглядного решения задач на формулу полной вероятности и формулу Байеса удобно использовать \textbf{деревья вероятностей}.

\begin{example}[Диагностика спама]
Пусть $8\%$ сообщений являются спамом (гипотеза $H_1$), а $92\%$ — нормальными письмами (гипотеза $H_2$). Вероятность того, что фильтр пометит спам как спам (событие $A$), равна $0{,}95$. Вероятность ложного срабатывания на нормальном письме равна $0{,}01$. 

Найдем общую вероятность срабатывания фильтра $P(A)$ и вероятность того, что помеченное письмо — действительно спам $P(H_1|A)$.

\begin{center}
\begin{tikzpicture}[
  level 1/.style={sibling distance=6cm, level distance=2.5cm},
  level 2/.style={sibling distance=3cm, level distance=2.5cm},
  every node/.style={align=center, fill=white, inner sep=2pt}
]
\node {Поступающее\\письмо}
  child { node {$H_1$ (Спам)}
    child { node {$A$ (Метка)} edge from parent node[left=2pt] {$P(A|H_1) = 0{,}95$} }
    child { node {$\overline{A}$ (Нет)} edge from parent node[right=2pt] {$0{,}05$} }
    edge from parent node[left=5pt, fill=none] {$P(H_1) = 0{,}08$}
  }
  child { node {$H_2$ (Норма)}
    child { node {$A$ (Метка)} edge from parent node[left=2pt] {$P(A|H_2) = 0{,}01$} }
    child { node {$\overline{A}$ (Нет)} edge from parent node[right=2pt] {$0{,}99$} }
    edge from parent node[right=5pt, fill=none] {$P(H_2) = 0{,}92$}
  };
\end{tikzpicture}
\end{center}

По формуле полной вероятности суммируем все пути, ведущие к событию $A$:
$$P(A) = 0{,}08 \cdot 0{,}95 + 0{,}92 \cdot 0{,}01 = 0{,}076 + 0{,}0092 = 0{,}0852.$$

По формуле Байеса вычисляем долю «истинного спама» среди всех сработавших меток:
$$P(H_1|A) = \frac{P(H_1)P(A|H_1)}{P(A)} = \frac{0{,}076}{0{,}0852} \approx 0{,}892 \quad (89{,}2\%).$$
\end{example}